\documentclass[../../syllabus_1.tex]{subfiles}
\begin{document}
\chapter{Proportionnalité}

\section{Tableaux de proportionnalite}

Un petit vendeur de pâtisseries vend ses croissants à \(1,5\)€.

\begin{enumerate}[itemsep=1.5cm]
  \item Combien vas-tu payer pour \(6\) croissants~?
  \item Complète le tableau suivant.
  \begin{center}
    \begin{tblr}{
      hlines,
      vlines,
      columns={1.5cm, c},
      rows={1cm,m},
      column{1}={4cm, gray!40, font=\bfseries},
  }
  Nombre de croissants & 0 & 1 & & 8 & & \\ 
  Prix (en €) & & & 6 & & 27 & 54 \\
\end{tblr}

  \end{center}
\end{enumerate}

\dotsforanswer{3}[Deux grandeurs directement proportionnelles sont]

\begin{exercise}
  Vrai ou faux~?
  \begin{enumerate}[itemsep=.4cm]
    \item Le nombre de bouteilles d'eau achetées et le prix payé sont proportionnels \dotfill
    \item La taille d'une personne est proportionnelle à son âge \dotfill
    \item Le salaire d'un professeur et le nombre de ses élèves en classe sont proportionnels \dotfill
    \item Le prix d'un livre et le nombre de pages contenues dans ce livre sont proportionnels \dotfill
    \item Le numéro d'un épisode de ta série préférée est proportionnel au temps de cet épisode \dotfill
    \item Le périmètre d'un carré est proportionnel à la mesure de son côté \dotfill
  \end{enumerate}
\end{exercise}

\begin{exercise}
  En t'aidant de la définition des grandeurs proportionnelles, complète les tableaux ci-dessous en indiquant sur les flèches les produits effectués.
  \begin{center}
    \includegraphics{chapitres/proportionnalite/img/tableauxprop.png}
  \end{center}
\end{exercise}

\dotsforanswer{2}[Le coefficient de proportionnalité (\(k\)) entre deux grandeurs proportionnelles (\(x\) et \(y\)) est]

\exemples
\begin{multicols}{2}
    \begin{center}
      \begin{tblr}{
        hlines,
        vlines,
        columns={1cm, c},
        rows={1cm,m},
        column{1}={gray!40, font=\bfseries},
    }
    \(x\) & 3 & 5 & 7 \\
    \(y\) & 12 & 20 & 28 \\
\end{tblr}
  \end{center}
  
    \begin{center}
      \begin{tblr}{
        hlines,
        vlines,
        columns={1cm, c},
        rows={1cm,m},
        column{1}={gray!40, font=\bfseries},
    }
    \(x\) & 2 & 8 & 9 \\
    \(y\) & 5 & 20 & 22,5  \\
\end{tblr}
  \end{center}
\end{multicols}

\begin{exercise}
  Détermine le coefficient de proportionnalité reliant le \(2\) grandeurs de chaque tableau de proportionnalité ci-dessous et utilise-le pour compléter les cases vides.
  \begin{center}
    \includegraphics{chapitres/proportionnalite/img/tableauxprop2.png}
  \end{center}
\end{exercise}

\newpage
\begin{exercise}
  Dans chaque cas, détermine si les grandeurs \(x\) et \(y\) sont proportionnelles. Si oui, détermine le coefficient de proportionnalité. 
  \begin{center}
    \includegraphics{chapitres/proportionnalite/img/tableauxprop3.png}
  \end{center}
\end{exercise}

\begin{exercise}
  Complète les tableaux suivants grâce au coefficient de proportionnalité.
  \begin{center}
    \includegraphics{chapitres/proportionnalite/img/tableauxprop4.png}
  \end{center}
\end{exercise}

\newpage
\section{Règle de trois}
\begin{exercise}
  On sait que \(4\) kilos de pommes coûtent \(6€\). Que coûtent \(7\) kilos de pommes~?\\ 
  Résous en complétant le tableau ci-dessous.

  \begin{center}
    \begin{tblr}{
            hlines,
            vlines,
            columns={3cm, c},
            rows={1cm,m},
            row{1}={gray!40, font=\bfseries},
        }
        Masse (\(kg\)) & Prix (€)\\
        4 & 6 \\
        1 & \\
        7 & \\
    \end{tblr}
\end{center}
\end{exercise}

\vfil

\begin{exercise}
  Ma voiture consomme, en moyenne, \(4\) litres de carburant pour \(100km\). Quelle distance pourrais-je parcourir avec un plein de \(50\) litres~?
  
\end{exercise}

\vfil

\begin{exercise}
  Pour une exposition, un groupe de \(12\) élèves a payé \(84€\). Qu'auraient payé \(17\) élèves~?
\end{exercise}

\newpage
\section{Graphiques et proportionnalité}
\begin{exercise}
  Voici les tarifs d'entrée à la piscine de la commune.
  \begin{center}
    \begin{tblr}{
            hlines,
            vlines,
            columns={1cm, c},
            rows={1cm,m},
            column{1}={3cm, gray!40, font=\bfseries},
        }
        Nombre d'entrées & 1 & 2 & 3 & 4 & 5 & 6 & 7 & 8 & 9 & 10 \\
        Prix entrée enfants & 3 & 6 & 9 & 12 & 15 & 18 & 21 & 24 & 27 & 30 \\
    \end{tblr}
\end{center}

\begin{enumerate}
  \item Est-ce un tableau de proportionnalité? Justifie. 
  \item Construis le graphique évolutif qyu traduit la situation. 
\end{enumerate}
\begin{center}
  \begin{tikzpicture}[scale=0.25]
      % Axes
      \draw[->] (0,-0.20) -- (0,30) node[above] {$Prix$ (en €)};
      \draw[->] (-0.20,0) -- (50,0) node[right] {$Nombre d'entrées$};
      % Grille
      \draw[gray, dotted] (0,0) grid (50,30);
      % Graduations
      \foreach \x in {1, 2,3,4,5,6,7,8,9,10} %\TODO: graduation
          \draw (\x,-0.1) -- (\x,0.1) node[below=2pt] {\x};
      \foreach \y in {5,10,15,20,25,30}
          \draw (-0.1,\y) -- (0.1,\y) node[left=2pt] {\y};
  \end{tikzpicture}
\end{center}
\end{exercise}

%  Voici les tarifs d'entrée à la piscine de la commune.
%\end{enumerate}
%
%\begin{longtable}[]{@{}
%  >{\raggedright\arraybackslash}p{(\columnwidth - 20\tabcolsep) * \real{0.3094}}
%  >{\raggedright\arraybackslash}p{(\columnwidth - 20\tabcolsep) * \real{0.0670}}
%  >{\raggedright\arraybackslash}p{(\columnwidth - 20\tabcolsep) * \real{0.0670}}
%  >{\raggedright\arraybackslash}p{(\columnwidth - 20\tabcolsep) * \real{0.0670}}
%  >{\raggedright\arraybackslash}p{(\columnwidth - 20\tabcolsep) * \real{0.0711}}
%  >{\raggedright\arraybackslash}p{(\columnwidth - 20\tabcolsep) * \real{0.0711}}
%  >{\raggedright\arraybackslash}p{(\columnwidth - 20\tabcolsep) * \real{0.0711}}
%  >{\raggedright\arraybackslash}p{(\columnwidth - 20\tabcolsep) * \real{0.0711}}
%  >{\raggedright\arraybackslash}p{(\columnwidth - 20\tabcolsep) * \real{0.0711}}
%  >{\raggedright\arraybackslash}p{(\columnwidth - 20\tabcolsep) * \real{0.0711}}
%  >{\raggedright\arraybackslash}p{(\columnwidth - 20\tabcolsep) * \real{0.0632}}@{}}
%\toprule()
%\begin{minipage}[b]{\linewidth}\raggedright
%Nombre d'entrées
%\end{minipage} & \begin{minipage}[b]{\linewidth}\raggedright
%\[1\]
%\end{minipage} & \begin{minipage}[b]{\linewidth}\raggedright
%\[2\]
%\end{minipage} & \begin{minipage}[b]{\linewidth}\raggedright
%\[3\]
%\end{minipage} & \begin{minipage}[b]{\linewidth}\raggedright
%\[4\]
%\end{minipage} & \begin{minipage}[b]{\linewidth}\raggedright
%\[5\]
%\end{minipage} & \begin{minipage}[b]{\linewidth}\raggedright
%\[6\]
%\end{minipage} & \begin{minipage}[b]{\linewidth}\raggedright
%\[7\]
%\end{minipage} & \begin{minipage}[b]{\linewidth}\raggedright
%\[8\]
%\end{minipage} & \begin{minipage}[b]{\linewidth}\raggedright
%\[9\]
%\end{minipage} & \begin{minipage}[b]{\linewidth}\raggedright
%\[10\]
%\end{minipage} \\
%\midrule()
%\endhead
%Prix entrée enfants & \(3\) & \(6\) & \(9\) & \(12\) & \(15\) & \(18\) &
%\(21\) & \(24\) & \(27\) & \(30\) \\
%\bottomrule()
%\end{longtable}
%
%\begin{enumerate}
%\def\labelenumi{\alph{enumi}.}
%\item
%  Est-ce un tableau de proportionnalité~? Justifie.
%\item
%  Construis le graphique évolutif qui traduit la situation.
%\end{enumerate}
%
%\includegraphics[width=5.22545in,height=4.74208in]{media/image3.tmp}
%
%\begin{enumerate}
%\def\labelenumi{\arabic{enumi}.}
%\setcounter{enumi}{9}
%\item
%  Parmi les graphiques suivants, quel est celui qui traduit une
%  situation de proportionnalité~? Justifie en disant pourquoi les autres
%  ne le sont pas.
%\end{enumerate}
%
%\includegraphics[width=3.63333in,height=0.88333in]{media/image4.png}
%
%\begin{enumerate}
%\def\labelenumi{\arabic{enumi}.}
%\setcounter{enumi}{13}
%\item
%  L'Échelle
%\end{enumerate}
%
%L'échelle
%\ldots\ldots\ldots\ldots\ldots\ldots\ldots\ldots\ldots\ldots\ldots\ldots\ldots\ldots\ldots\ldots\ldots\ldots\ldots\ldots\ldots\ldots\ldots\ldots\ldots\ldots\ldots.\ldots\ldots\ldots\ldots\ldots\ldots\ldots.\ldots\ldots\ldots\ldots\ldots\ldots\ldots\ldots\ldots\ldots\ldots\ldots\ldots\ldots\ldots\ldots\ldots\ldots\ldots.
%
%\ldots\ldots\ldots\ldots\ldots\ldots\ldots\ldots\ldots\ldots\ldots\ldots\ldots\ldots\ldots\ldots\ldots\ldots\ldots\ldots\ldots\ldots\ldots\ldots\ldots\ldots\ldots\ldots\ldots\ldots\ldots\ldots..\ldots\ldots\ldots\ldots\ldots\ldots\ldots\ldots.\ldots\ldots\ldots\ldots\ldots\ldots\ldots\ldots\ldots\ldots\ldots\ldots\ldots\ldots\ldots\ldots\ldots\ldots.
%
%\[Echelle = \frac{longueur\ représentée}{longueur\ réel}\]
%
%\includegraphics[width=2.77014in,height=2.30278in]{media/image5.png}
%
%\begin{longtable}[]{@{}
%  >{\raggedright\arraybackslash}p{(\columnwidth - 2\tabcolsep) * \real{0.5000}}
%  >{\raggedright\arraybackslash}p{(\columnwidth - 2\tabcolsep) * \real{0.5000}}@{}}
%\toprule()
%\begin{minipage}[b]{\linewidth}\raggedright
%\textbf{Mesure sur la carte (cm)}
%\end{minipage} & \begin{minipage}[b]{\linewidth}\raggedright
%\textbf{Distance réelle (cm)}
%\end{minipage} \\
%\midrule()
%\endhead
%1 & \\
%2 & \\
%3 & \\
%\bottomrule()
%\end{longtable}
%
%Exemple~: La Belgique est représenté à l'échelle~: 1 / 4 000 000
%
%L'échelle est
%\ldots\ldots\ldots\ldots\ldots\ldots\ldots\ldots\ldots\ldots\ldots\ldots\ldots\ldots\ldots\ldots\ldots\ldots\ldots\ldots\ldots\ldots\ldots\ldots\ldots\ldots\ldots\ldots\ldots\ldots\ldots\ldots\ldots\ldots\ldots\ldots\ldots\ldots\ldots\ldots\ldots\ldots\ldots\ldots\ldots\ldots\ldots..
%
%\begin{enumerate}
%\def\labelenumi{\arabic{enumi}.}
%\setcounter{enumi}{14}
%\item
%  Exercices mélangés
%\end{enumerate}
%
%\begin{enumerate}
%\def\labelenumi{\arabic{enumi}.}
%\item
%  Sur un plan dont l'échelle est \(\frac{1}{2000}\), un champ
%  rectangulaire est représenté par un rectangle mesurant \(2cm\) de long
%  et \(1,4cm\) de large. Quelles sont les dimensions réelles de ce
%  champ, en mètres~? Quelle est l'aire de ce champ en mètres carrés~?
%\end{enumerate}
%
%\begin{enumerate}
%\def\labelenumi{\arabic{enumi}.}
%\setcounter{enumi}{1}
%\item
%  En 2 heures, Pierre a parcouru 180 km. À vitesse constante, combien de
%  kilomètres parcourra-t-il en 5 heures~?
%\item
%  Un peintre a peint 3pièces de même superficie en 6 jours. À vitesse
%  constante, de combien de jours a-t-il besoin pour peindre 4 pièces
%  identiques?
%\item
%  Pour imprimer 28 journaux de l\textquotesingle école, li a fallu 112
%  feuilles. Combien faudra-t-il de feuilles pour imprimer 84 journaux?
%\item
%  Anne roule à vélo à 15 km/h. Àvitesse constante, combien de temps
%  mettra-t-elle pour parcourir 60 km ?
%\end{enumerate}
%
%\begin{enumerate}
%\def\labelenumi{\arabic{enumi}.}
%\setcounter{enumi}{5}
%\item
%  Chez Colruyt, un lot de trois bouteilles de vin est affiché à \(21€\).
%  On peut acheter les bouteilles à la pièce. Combien vais-je payer pour
%  l'achat de \(10\) bouteilles~?
%\item
%  Exercice issu d'un CE1D.
%\end{enumerate}
%
%Tableau A
%
%\begin{longtable}[]{@{}
%  >{\raggedright\arraybackslash}p{(\columnwidth - 2\tabcolsep) * \real{0.5000}}
%  >{\raggedright\arraybackslash}p{(\columnwidth - 2\tabcolsep) * \real{0.5000}}@{}}
%\toprule()
%\begin{minipage}[b]{\linewidth}\raggedright
%\[x\]
%\end{minipage} & \begin{minipage}[b]{\linewidth}\raggedright
%\[y\]
%\end{minipage} \\
%\midrule()
%\endhead
%\(3\) & \(9\) \\
%\(2,5\) & \(7,5\) \\
%\(9\) & \(27\) \\
%\(10,1\) & \(30,3\) \\
%\bottomrule()
%\end{longtable}
%
%Tableau B
%
%\begin{longtable}[]{@{}
%  >{\raggedright\arraybackslash}p{(\columnwidth - 2\tabcolsep) * \real{0.5000}}
%  >{\raggedright\arraybackslash}p{(\columnwidth - 2\tabcolsep) * \real{0.5000}}@{}}
%\toprule()
%\begin{minipage}[b]{\linewidth}\raggedright
%\[x\]
%\end{minipage} & \begin{minipage}[b]{\linewidth}\raggedright
%\[y\]
%\end{minipage} \\
%\midrule()
%\endhead
%\(1\) & \(3\) \\
%\(5\) & \(7\) \\
%\(17\) & \(19\) \\
%\(35\) & \(37\) \\
%\bottomrule()
%\end{longtable}
%
%\textbf{COCHE} la case du tableau qui montre une proportionnalité
%directe entre la grandeur \(x\) et la grandeur \(y\).
%
%Pour ce tableau, \textbf{ÉCRIS} le coefficient de proportionnalité~:
%
%\begin{enumerate}
%\def\labelenumi{\arabic{enumi}.}
%\setcounter{enumi}{6}
%\item
%  Vrai ou Faux~? Les situations suivantes sont-elles des situations de
%  proportionnalité~?
%\end{enumerate}
%
%\begin{longtable}[]{@{}
%  >{\raggedright\arraybackslash}p{(\columnwidth - 4\tabcolsep) * \real{0.7980}}
%  >{\raggedright\arraybackslash}p{(\columnwidth - 4\tabcolsep) * \real{0.1037}}
%  >{\raggedright\arraybackslash}p{(\columnwidth - 4\tabcolsep) * \real{0.0983}}@{}}
%\toprule()
%\begin{minipage}[b]{\linewidth}\raggedright
%\end{minipage} & \begin{minipage}[b]{\linewidth}\raggedright
%Vrai
%\end{minipage} & \begin{minipage}[b]{\linewidth}\raggedright
%Faux
%\end{minipage} \\
%\midrule()
%\endhead
%Le gain au lotto par rapport à la mise de départ. & & \\
%La masse d\textquotesingle une personne par rapport à son âge. & & \\
%Le prix des oranges par rapport au nombre de kg achetés. & & \\
%La distance parcourue par un automobiliste par rapport à sa vitesse de
%
%conduite (en supposant qu\textquotesingle il roule à une vitesse
%constante de 70km/h). & & \\
%La quantité de farine utilisée pour réaliser la recette de la pâte à
%crêpes
%
%par rapport au nombre de crêpes voulu. & & \\
%Ta movenne au cours de mathématique par rapport au temps de travail
%
%à la maison. & & \\
%Le nombre de paniers de basket marqués par rapport au temps
%d\textquotesingle une partie. & & \\
%Le salaire d\textquotesingle un professeur par rapport au taux de
%réussite d\textquotesingle élèves à son cours. & & \\
%\bottomrule()
%\end{longtable}
%
%\begin{enumerate}
%\def\labelenumi{\arabic{enumi}.}
%\setcounter{enumi}{7}
%\item
%  Pour télécharger \(3\) chansons sur internet, à une certaine époque,
%  il fallait \(1\) minute.
%\end{enumerate}
%
%\begin{enumerate}
%\def\labelenumi{\alph{enumi}.}
%\item
%  Complète, en te basant sur ce temps moyen de téléchargement, le
%  tableau de proportionnalité suivant~:
%\end{enumerate}
%
%\begin{longtable}[]{@{}
%  >{\raggedright\arraybackslash}p{(\columnwidth - 2\tabcolsep) * \real{0.4966}}
%  >{\raggedright\arraybackslash}p{(\columnwidth - 2\tabcolsep) * \real{0.5034}}@{}}
%\toprule()
%\begin{minipage}[b]{\linewidth}\raggedright
%Nombre de chansons
%\end{minipage} & \begin{minipage}[b]{\linewidth}\raggedright
%Durée de téléchargement
%
%(en secondes)
%\end{minipage} \\
%\midrule()
%\endhead
%& 120 \\
%\(9\) & \\
%& 500 \\
%& 60 \\
%1 & \\
%\bottomrule()
%\end{longtable}
%
%\begin{enumerate}
%\def\labelenumi{\alph{enumi}.}
%\setcounter{enumi}{1}
%\item
%  Calcule le nombre de chansons que tu pourrais télécharger, à la même
%  vitesse, en une demi-heure.
%\end{enumerate}
%
%\begin{enumerate}
%\def\labelenumi{\arabic{enumi}.}
%\setcounter{enumi}{8}
%\item
%  Mathis achète la maquette d\textquotesingle une voiture reproduite à
%  l\textquotesingle échelle 1:22.
%\end{enumerate}
%
%\begin{enumerate}
%\def\labelenumi{\alph{enumi})}
%\item
%  Le modèle réduit mesure après construction 19,8 cm. Quelle est la
%  longueur réelle de cette voiture ?
%\item
%  Si la taille réelle d\textquotesingle une voiture était de 5,17 m,
%  quelle serait la longueur de la maquette représentant cette voiture en
%  utilisant la m ê m e échelle ?
%\end{enumerate}
%
%\begin{enumerate}
%\def\labelenumi{\arabic{enumi}.}
%\setcounter{enumi}{9}
%\item
%  Si le monde était un village\ldots{}
%\end{enumerate}
%
%\begin{quote}
%En 1990, Donella Meadows écrivait un texte intitulé «~State of the
%Village Report~», traduit en français sous le nom «~Si le monde était un
%village~». Il ramenait toutes des statistiques du monde entier à une
%population plus petite, celle d'un village de \textbf{100 personnes},
%tout en conservant les proportions de tous les peuples existants sur la
%terre. Il a été retravaillé plusieurs fois au cours des ans pour que les
%chiffres collent à la réalité de la population grandissante, et nous
%pouvons ainsi avoir une vision plus rapide de ces statistiques. Par
%exemple, si le monde était un village de 100 personnes, 48 seraient des
%hommes et 52 des femmes. Ou encore, si le monde était un village de 100
%personnes, 6 personnes possèderaient 59\% de la richesse totale et tous
%les 6 seraient originaires des USA.
%
%Voici quelques chiffres de ce texte. Réponds aux questions qui suivent
%en laissant ta démarche visible.
%\end{quote}
%
%\begin{enumerate}
%\def\labelenumi{\alph{enumi}.}
%\item
%  \emph{Si le monde était un village de 100 personnes, 68 respirent un
%  air pur et 32 un air pollué.} Que seraient ces chiffres si le monde
%  était un village de 25 personnes~?
%\item
%  \emph{Si le monde était un village de 100 personnes, 83 personnes ont
%  accès à l'eau potable alors que 17 n'ont pas accès à de l'eau
%  potable.} Combien de personnes auraient accès à l'eau potable si le
%  monde était constitué de 40 personnes~? Arrondis.
%\item
%  \emph{Si le monde était un village de 100 personnes, 48 personnes ne
%  peuvent pas parler ou agir en âme et conscience (prison, torture ou
%  mort à la clef), alors que 52 personnes le peuvent.} Que seraient ces
%  nombres pour une population mondiale de 350 personnes~?
%\item
%  \emph{Si le monde était un village de 100 personnes, 7 personnes ont
%  un ordinateur, mais 93 personnes n'en ont pas.} Combien de personnes
%  auraient accès à un ordinateur si le monde était constitué de 117
%  personnes~? Arrondis.
%\item
%  \emph{Si le monde était un village de 100 personnes, 99 personnes
%  n'auraient pas de formation universitaire, alors qu'une seule personne
%  aurait un diplôme universitaire}. Quel pourcentage de la population
%  mondiale a suivi un parcours universitaire~?
%\end{enumerate}
%
%\begin{enumerate}
%\def\labelenumi{\arabic{enumi}.}
%\setcounter{enumi}{10}
%\item
%  Un robinet ouvert à fond remplit un seau de \(12\) litres en \(18\)
%  secondes. Etudie le temps qu'il faudra pour remplir des seaux de
%  différentes capacités.
%\end{enumerate}
%
%\begin{enumerate}
%\def\labelenumi{\alph{enumi}.}
%\item
%  Complète le tableau suivant correspondant à la situation donnée.
%\end{enumerate}
%
%\begin{longtable}[]{@{}
%  >{\raggedright\arraybackslash}p{(\columnwidth - 10\tabcolsep) * \real{0.2997}}
%  >{\raggedright\arraybackslash}p{(\columnwidth - 10\tabcolsep) * \real{0.1400}}
%  >{\raggedright\arraybackslash}p{(\columnwidth - 10\tabcolsep) * \real{0.1400}}
%  >{\raggedright\arraybackslash}p{(\columnwidth - 10\tabcolsep) * \real{0.1401}}
%  >{\raggedright\arraybackslash}p{(\columnwidth - 10\tabcolsep) * \real{0.1400}}
%  >{\raggedright\arraybackslash}p{(\columnwidth - 10\tabcolsep) * \real{0.1401}}@{}}
%\toprule()
%\begin{minipage}[b]{\linewidth}\raggedright
%Quantité d'eau (en litres)
%\end{minipage} & \begin{minipage}[b]{\linewidth}\raggedright
%\[3\]
%\end{minipage} & \begin{minipage}[b]{\linewidth}\raggedright
%\[5\]
%\end{minipage} & \begin{minipage}[b]{\linewidth}\raggedright
%\[6\]
%\end{minipage} & \begin{minipage}[b]{\linewidth}\raggedright
%\[12\]
%\end{minipage} & \begin{minipage}[b]{\linewidth}\raggedright
%\[20\]
%\end{minipage} \\
%\midrule()
%\endhead
%Temps (en secondes) & & & & \(18\) & \\
%\bottomrule()
%\end{longtable}
%
%\begin{enumerate}
%\def\labelenumi{\alph{enumi}.}
%\setcounter{enumi}{1}
%\item
%  Trace le graphique liant la quantité d'eau déversée par le robinet en
%  fonction du temps.
%\end{enumerate}
%
%\begin{quote}
%\includegraphics[width=4.63333in,height=4.475in]{media/image6.png}
%\end{quote}
%
%\begin{enumerate}
%\def\labelenumi{\alph{enumi}.}
%\setcounter{enumi}{2}
%\item
%  Justifie à l'aide du tableau \textbf{et} du graphique qu'il s'agit
%  d'une situation de proportionnalité.
%\end{enumerate}
%
%\begin{enumerate}
%\def\labelenumi{\arabic{enumi}.}
%\setcounter{enumi}{11}
%\item
%  Complète le portrait du Minion (page suivante) en calculant les
%  valeurs ci-dessous.
%\end{enumerate}
%
%\begin{longtable}[]{@{}
%  >{\raggedright\arraybackslash}p{(\columnwidth - 8\tabcolsep) * \real{0.2133}}
%  >{\raggedright\arraybackslash}p{(\columnwidth - 8\tabcolsep) * \real{0.1966}}
%  >{\raggedright\arraybackslash}p{(\columnwidth - 8\tabcolsep) * \real{0.1967}}
%  >{\raggedright\arraybackslash}p{(\columnwidth - 8\tabcolsep) * \real{0.1966}}
%  >{\raggedright\arraybackslash}p{(\columnwidth - 8\tabcolsep) * \real{0.1967}}@{}}
%\toprule()
%\begin{minipage}[b]{\linewidth}\raggedright
%Une boite de bonbons coûte \(1,99€\). Trois boites coûtent
%\end{minipage} & \begin{minipage}[b]{\linewidth}\raggedright
%3,87
%
%\includegraphics[width=1.04176in,height=0.64172in]{media/image7.tmp}
%\end{minipage} & \begin{minipage}[b]{\linewidth}\raggedright
%4,90€
%
%\includegraphics[width=0.71673in,height=0.68339in]{media/image8.tmp}
%\end{minipage} & \begin{minipage}[b]{\linewidth}\raggedright
%5,97€
%
%\includegraphics[width=0.93341in,height=0.60839in]{media/image9.tmp}
%\end{minipage} & \begin{minipage}[b]{\linewidth}\raggedright
%1,10€
%
%\includegraphics[width=1.05009in,height=0.47504in]{media/image10.tmp}
%\end{minipage} \\
%\midrule()
%\endhead
%Je prends mon pouls. En \(20\)s, je compte 25 pulsations. Si mon pouls
%est régulier, en une minute, je compte & 65 pulsations
%
%\includegraphics[width=1.36678in,height=0.61672in]{media/image11.tmp} &
%75 pulsations
%
%\includegraphics[width=1.27511in,height=0.62505in]{media/image12.tmp} &
%125 pulsations
%
%\includegraphics[width=1.21677in,height=0.71673in]{media/image13.tmp} &
%105 pulsations
%
%\includegraphics[width=1.29178in,height=0.63339in]{media/image14.tmp} \\
%Sur l'autoroute, à vitesse constante, je parcours 6km en 3min. Donc en
%5min, je parcours & 10km
%
%\includegraphics[width=1.1001in,height=0.55838in]{media/image15.tmp} &
%8km
%
%\includegraphics[width=1.03342in,height=0.52505in]{media/image16.tmp} &
%2km
%
%\includegraphics[width=1.14177in,height=0.48337in]{media/image17.tmp} &
%11km
%
%\includegraphics[width=1.14177in,height=0.53338in]{media/image18.tmp} \\
%10 gâteaux identiques pèsent 220g, donc 15 gâteaux pèsent &
%440g\includegraphics[width=1.10843in,height=0.56672in]{media/image19.tmp}
%& 225g
%
%\includegraphics[width=0.98342in,height=0.50838in]{media/image20.tmp} &
%330g
%
%\includegraphics[width=0.96675in,height=0.7584in]{media/image21.tmp} &
%235g
%
%\includegraphics[width=1.30845in,height=0.51671in]{media/image22.tmp} \\
%Le son parcourt 300m en 1s. En une minute, il parcourt & 1800m
%
%\includegraphics[width=1.2001in,height=0.94175in]{media/image23.tmp} &
%30km
%
%\includegraphics[width=1.24177in,height=0.90841in]{media/image24.tmp} &
%359m
%
%\includegraphics[width=1.21677in,height=0.67506in]{media/image25.tmp} &
%18km
%
%\includegraphics[width=1.18333in,height=1.19167in]{media/image26.tmp} \\
%Pour faire un gâteau pour 4 personnes, il faut 100g de chocolat. Pour 6
%personnes, il en faudra & 200g
%
%\includegraphics[width=1.2001in,height=1.07509in]{media/image27.tmp} &
%102g
%
%\includegraphics[width=1.14177in,height=1.05843in]{media/image28.tmp} &
%50g
%
%\includegraphics[width=1.175in,height=1.05in]{media/image29.tmp} & 150g
%
%\includegraphics[width=1.25844in,height=1.06676in]{media/image30.tmp} \\
%\bottomrule()
%\end{longtable}
%
%Dessin~:
%
%\includegraphics[width=3.74167in,height=5.13333in]{media/image31.tmp}
\end{document}